Now
@@M0@@

Moreover, @@M1@@; we shall denote by @@M2@@ the class in @@M3@@ of @@M4@@. Relation (1.11.1) therefore finally gives:
\sgaStableTarget{sga6.e06.eq.1-11-2}@@M5@@

We then consider the complex:
@@M6@@
whose arrows are defined by:
@@M7@@
If we regard @@M8@@ as a graded @@M9@@-Module concentrated in degree @@M10@@, this is a complex of graded @@M11@@-Modules, the morphisms themselves being graded.

This complex is a resolution of @@M12@@: on an open set where @@M13@@ is free, @@M14@@ identifies with a polynomial Algebra @@M15@@, and the complex with the Koszul complex defined by the @@M16@@-regular sequence @@M17@@; we know that the latter is a resolution of @@M18@@.

Writing that each homogeneous component of the complex of degree @@M19@@ is acyclic, we therefore obtain relations in @@M20@@:
\sgaStableTarget{sga6.e06.eq.1-11-3}@@M21@@
(putting @@M22@@ for @@M23@@).

Suppose relation (1.11.2) holds for @@M24@@; we shall show that it still holds for @@M25@@, which will allow us to descend by induction to @@M26@@ (we shall see later that this suffices to imply the vanishing of the @@M27@@).

